(a) Only finitely many points of $X$ are singular, and $\overline{\mathcal O}_P=\mathcal O_P$ elsewhere. A tuple in $\bigoplus_P\overline{\mathcal O}_P^*/\mathcal O_P^*$ determines a Cartier divisor by using representatives as local equations near the singular points and $1$ away from them, shrinking the neighborhoods so the equations are units on the overlaps. Its pullback to $\overline X$ is zero. Changing a representative by a unit does not change the Cartier divisor.

If this divisor is principal, say $(f)$, then $(f)_{\overline X}=0$. Since $\overline X$ is an integral projective curve, $f\in k^*$, so the original tuple is zero. Thus the induced map to $\operatorname{Pic}X$ is injective. Conversely, if $\pi^*[D]=0$, subtract the principal divisor trivializing $\pi^*D$. Every local equation of the resulting divisor is a unit on the normalization, and its classes at the singular points give the required tuple. This proves exactness in the middle.

To prove surjectivity, represent a line bundle on $\overline X$ by a divisor. By the Chinese remainder theorem in an affine neighborhood of the finitely many points above the singular locus, choose a rational function having the same valuations at those points. Subtracting its divisor leaves a divisor supported where $\pi$ is an isomorphism, hence a Cartier divisor on $X$. Therefore
\[0\to\bigoplus_{P\in X}\overline{\mathcal O}_P^*/\mathcal O_P^*\to\operatorname{Pic}X\xrightarrow{\pi^*}\operatorname{Pic}\overline X\to0.\]

(b) Both normalizations are $\mathbb P^1$, whose Picard group is $\mathbb Z$. At a cusp, with parameter $t$, the local rings are $B=k[t]_{(t)}$ and $A=k[t^2,t^3]_{(t^2,t^3)}=k+t^2B$. The map $B^*\to k$, $u\mapsto u'(0)/u(0)$, is a surjective homomorphism to the additive group, with kernel $A^*$. Thus the first kernel in (a) is $\mathbb G_a$.

At a node, write the two points of the normalization as $t=a,b$. The same local-ring calculation as in \exref{II.6.7} gives $A^*=\{u\in B^*:u(a)=u(b)\}$, where $B$ is the semilocal normalization. The ratio $u(a)/u(b)$ induces $B^*/A^*\cong k^*$. Consequently the two exact sequences are $0\to\mathbb G_a\to\operatorname{Pic}X\to\mathbb Z\to0$ and $0\to\mathbb G_m\to\operatorname{Pic}X\to\mathbb Z\to0$, respectively.
